%% threats_to_validity.tex -- submission version

\section{Threats to Validity}
\label{sec:threats-to-validity}

We organize the limitations following the internal, external, construct,
conclusion, and reliability categories of Wohlin et al.~\cite{wohlin2012}.
The evidence hierarchy is explicit: native-framework executions are the
load-bearing RQ2 evidence; author-constructed recovery nodes and the
structural reauthorization proposal are supplementary.

\subsection{Internal Validity}

\paragraph{Native execution and matched controls.}
The canonical Reflexion matrix executes the framework's retry path at a fixed
commit and pairs every attack arm with a matched benign arm.  AutoGen 0.4.7
and Aider 0.86.2 likewise execute their native error/recovery flows and retain
path-completion indicators.  These controls reduce the risk that ordinary
code generation or tool use is misclassified as an attack, but they do not
remove model stochasticity or provider-side variation.

\paragraph{Author-constructed screens.}
The provenance/role/position contrast screen and the defense-aware screen use
an author-constructed recovery node.  They are mechanism and screening
experiments, not evidence that a named framework ships the exact vulnerable
node.  The defense-aware contrast is \DefenseUndefendedASR{} versus
\DefenseDefendedASR{} with Fisher's exact $p=\DefenseFisherP{}$; it is
therefore not presented as a statistically established
defense effect.  LangGraph recovery-node results are treated in the same way
because LangGraph does not ship that recovery primitive by default.

\paragraph{Measurement proxy.}
Attack success is classified from the framework artifact or model response
using a conservative, rule-based canary detector.  A response that merely
mentions a tool without expressing the canary behavior is not counted as a
success.  The native Aider and Reflexion logs provide a stronger ecological
check than a prompt-only reconstruction, while the remaining screens should
be interpreted as behavioral proxies rather than measurements of real-world
side effects.

\subsection{External Validity}

\paragraph{Framework sample.}
The source analysis covers six open-source frameworks selected for visible
retry, reflection, or error-recovery code.  This purposive sample supports a
comparative qualitative claim, not a prevalence estimate for all agent
frameworks.  Closed-source systems and frameworks with different recovery
architectures may behave differently.

\paragraph{Model and provider variation.}
The canonical matrix contains five provider aliases and shows substantial
heterogeneity (ASR \ReflexionMinASR{}--\ReflexionMaxASR{}).  AutoGen and Aider use smaller $n=50$ per
payload/arm cells, and the defense-aware and legitimate-recovery experiments
are screening samples.  The results should therefore be reported by
framework, model, and payload, rather than as one pooled model-independent
rate.

\paragraph{Snapshot disclosure.}
The OpenCode Go catalog records aliases, ownership, and one common
\code{created} timestamp.  Because that timestamp is identical across the
catalog entries, we interpret it as catalog publication time, not an
immutable checkpoint identifier.  The service documentation specifies
\code{/responses} for GPT-5.6 Luna, whereas the historical GPT artifacts were
created through \code{/chat/completions}.  A current one-request probe was
blocked with HTTP 403/edge code 1010 on both routes, so it cannot establish
protocol equivalence.  Billing records, endpoint metadata, and local raw logs
prove that calls were made, but cannot prove the exact backend checkpoint.
Three public Hugging Face repositories are recorded as snapshot candidates,
but no provider alias-to-revision mapping is disclosed.  Exact snapshot
reproducibility is consequently an explicit limitation.

\subsection{Construct Validity}

\paragraph{Trampoline construct.}
The trust-escalation trampoline is operationalized as the conjunction of two
source-level criteria: absence of a provenance distinction and application of
a high-trust recovery frame.  Partial cases are retained rather than forced
into the binary label.  The code excerpts, fixed commits, written rubric, and
deterministic checks are preserved.  This is a single-coder qualitative
analysis; independent re-coding and formal disagreement adjudication remain
limitations rather than completed claims.

\paragraph{Payload coverage.}
Five payload variants are used in the canonical Reflexion matrix and native
framework runs.  They cover the same semantic directive with different
surface forms, but cannot represent the full space of adaptive or multilingual
attacks.  The defense-aware screen is single-round; a multi-round attacker
that learns from rejections could achieve a different rate.

\paragraph{Legitimate recovery.}
The legitimate-recovery screen contains \LegalRecoveryTasks{} deterministic
tasks across seven failure categories, with \LegalRecoverySuccesses{}
quality-passing recoveries and zero false positives.  It is reported as
screening evidence only.  The accompanying power analysis is prospective
planning information, not post-hoc confirmation.

\subsection{Conclusion Validity}

All reported intervals are Wilson 95\% intervals, and matched binary
comparisons use two-sided Fisher exact tests.  The contrast and defense-aware
screens are explicitly exploratory or descriptive where power is insufficient.
The power-analysis artifact records that the small provenance contrast and
the observed defense contrast would require substantially larger samples for
confirmatory inference.  We therefore avoid claims that provenance tagging
or the completed defense screen has a proven population-level effect.

\subsection{Reliability}

The replication package retains raw per-trial logs, canonical manifests,
endpoint metadata, artifact hashes, the GPT endpoint probe, the model-snapshot
audit, the power analysis, the provenance manifest, the source-code audit, and
the number-generation script.  Credentials are not written to these
artifacts.  All numerical tables in this submission are generated from the
JSON logs; any older draft section that cannot be traced to those logs is
excluded from the quantitative claims.

\paragraph{Structural-defense scope.}
The \reauthgate{} section specifies a useful authorization invariant, but this
paper does not evaluate a native-framework port with raw logs.  Claims of zero
ASR, framework-agnosticity, or negligible overhead are therefore design
hypotheses and future-work requirements, not findings used to support the main
RQ2/RQ3 conclusions.
